\chapter{Entropy and Data Compression}

\section{Probability Review}
This section will review some basic concepts of probability theory that were covered in the previous course CS-101.

\subsection{Definitions}
\begin{definition}[Sample Space]
    A sample space is a set $\Omega$ of possible outcomes of a random experiment. Each outcome $\omega \in \Omega$ is called an elementary event.
\end{definition}

\begin{definition}[Event]
    An event is a subset $E \subseteq \Omega$ of the sample space $\Omega$. The probability of an event $E$ is denoted:
    \[
        P(E) = \frac{|E|}{|\Omega|}
    \]
\end{definition}
Note that this definition of probability assumes that all outcomes in the sample space are equally likely, which is often the case in simple experiments.

\begin{eg}
    Consider the experiment of tossing a fair coin three times. The sample space is:
    \[
        \Omega = \{HHH, HHT, HTH, HTT, THH, THT, TTH, TTT\}
    \]
    Each outcome has a probability of $\frac{1}{8}$. An example of an event is:
    \[
        E = \{\text{at least two heads}\} = \{HHH, HHT, HTH, THH\}
    \]
    which has a probability of $P(E) = \frac{4}{8} = \frac{1}{2}$.
\end{eg}

\begin{definition}[Probability Mass Function]
    A probability mass function (PMF) is a function $P: \Omega \to [0,1]$ that assigns a probability to each outcome in the sample space, such that:
    \[
        \sum_{\omega \in \Omega} P(\omega) = 1
    \]
\end{definition}
Thus we can also define the probability of an event $E$ as:
\[
    P(E) = \sum_{\omega \in E} P(\omega)
\]

\subsection{Conditional Probability and Independence}
\begin{definition}[Conditional Probability]
    The conditional probability of an event $E$ given an event $F$ is defined as:
    \[
        P(E|F) = \frac{P(E \cap F)}{P(F)}
    \]
    provided that $P(F) > 0$.
\end{definition}
Is it often useful to see $F$ as the "new" sample space, and $E \cap F$ as the "new" event we are interested in. It can also be translated into english as: the probability of $E$ given that $F$ has occurred.

\begin{definition}[Independence]
    Two events $E$ and $F$ are said to be independent if:
    \[
        P(E \cap F) = P(E)P(F)
    \]
    Equivalently, if $P(F) > 0$, we can also say that $E$ and $F$ are independent if:
    \[
        P(E|F) = P(E)
    \]
\end{definition}

\begin{eg}
    Let's consider the experiment of tossing two fair coins. The sample space is:
    \[
        \Omega = \{HH, HT, TH, TT\}
    \]
    Let $E$ be the event "the first coin is heads" and $F$ be the event "the second coin is heads". We have:
    \[
        P(E) = \frac{2}{4} = \frac{1}{2}, \quad P(F) = \frac{2}{4} = \frac{1}{2}
    \]
    and
    \[
        P(E \cap F) = P(\{HH\}) = \frac{1}{4}
    \]
    Since $P(E)P(F) = \frac{1}{2} \cdot \frac{1}{2} = \frac{1}{4}$, we see that $P(E \cap F) = P(E)P(F)$, so $E$ and $F$ are independent.
\end{eg}

\begin{eg}
    Consider the experiment of rolling a fair six-sided die. Let $E$ be the event "the outcome is even" and $F$ be the event "the outcome is greater than 3". We have:
    \[
        P(E) = \frac{3}{6} = \frac{1}{2}, \quad P(F) = \frac{3}{6} = \frac{1}{2}
    \]
    and
    \[
        P(E \cap F) = P(\{\text{4, 6}\}) = \frac{2}{6} = \frac{1}{3}
    \]
    Since $P(E)P(F) = \frac{1}{2} \cdot \frac{1}{2} = \frac{1}{4}$, we see that $P(E \cap F) \neq P(E)P(F)$, so $E$ and $F$ are not independent.
\end{eg}

\begin{definition}[Disjoint Events]
    Two events $E$ and $F$ are said to be disjoint (or mutually exclusive) if:
    \[
        E \cap F = \emptyset
    \]
    In this case, we have:
    \[
        P(E \cup F) = P(E) + P(F)
    \]
\end{definition}

\subsection{Law of Total Probability and Bayes' Theorem}
\begin{definition}[Partition of the Sample Space]
    A collection of events $\{B_1, B_2, \ldots, B_n\}$ is called a partition of the sample space $\Omega$ if:
    \begin{citemize}
        \item $B_i \cap B_j = \emptyset$ for all $i \neq j$ (the events are mutually disjoint).
        \item $\bigcup_{i=1}^n B_i = \Omega$ (the events cover the entire sample space).
    \end{citemize}
\end{definition}

\begin{theorem}[Law of Total Probability]
    Let $\{B_1, B_2, \ldots, B_n\}$ be a partition of the sample space $\Omega$, meaning that $B_i \cap B_j = \emptyset$ for $i \neq j$ and $\bigcup_{i=1}^n B_i = \Omega$. Then for any event $E$, we have:
    \[
        P(E) = \sum_{i=1}^n P(E|B_i)P(B_i)
    \]
\end{theorem}
\begin{proof}
    Since $\{B_1, B_2, \ldots, B_n\}$ is a partition of $\Omega$, we can express $E$ as the union of its intersections with the $B_i$:
    \[
        E = (E \cap B_1) \cup (E \cap B_2) \cup \cdots \cup (E \cap B_n)
    \]
    Since the $B_i$ are mutually disjoint, the sets $E \cap B_i$ are also mutually disjoint. Therefore, we can write:
    \[
        P(E) = P\left( \bigcup_{i=1}^n (E \cap B_i) \right) = \sum_{i=1}^n P(E \cap B_i)
    \]
    Now, using the definition of conditional probability, we have:
    \[
        P(E|B_i) = \frac{P(E \cap B_i)}{P(B_i)} \quad \text{for } P(B_i) > 0
    \]
    Rearranging this gives:
    \[
        P(E \cap B_i) = P(E|B_i)P(B_i)
    \]
    Substituting this back into our expression for $P(E)$, we get:
    \[
        P(E) = \sum_{i=1}^n P(E|B_i)P(B_i)
    \]
\end{proof}

\begin{eg}
    Let's consider two factories that produce light bulbs.
    \begin{citemize}
        \item Factory A's bulbs work for over $5000$ hours in $99\%$ of the cases.
        \item Factory B's bulbs work for over $5000$ hours in $95\%$ of the cases.
        \item It is known that factory A produces $60\%$ of the bulbs.
    \end{citemize}
    Let's compute the probability that a randomly chosen bulb works for over $5000$ hours. We can use the law of total probability with the partition $\{B_A, B_B\}$, where $B_A$ is the event "the bulb is produced by factory A" and $B_B$ is the event "the bulb is produced by factory B". We have:
    \[
        P(\text{works}) = P(\text{works}|B_A)P(B_A) + P(\text{works}|B_B)P(B_B)
    \]
    Substituting the given probabilities, we get:
    \[
        P(\text{works}) = (0.99)(0.60) + (0.95)(0.40) = 0.594 + 0.38 = 0.974
    \]
\end{eg}
Sometimes we are interested in the reverse problem: given that a bulb works for over $5000$ hours, what is the probability that it was produced by factory A? This is where Bayes' theorem comes into play.

\begin{theorem}[Bayes' Theorem]
    Let $F, E \in \Omega$ be two events such that $P(E) > 0$. Then we have:
    \[
        P(F|E) = \frac{P(E|F)P(F)}{P(E)}
    \]
\end{theorem}
\begin{proof}
    By the definition of conditional probability, we have:
    \[
        P(F|E) = \frac{P(F \cap E)}{P(E)}
    \]
    Now, we can express $P(F \cap E)$ using the definition of conditional probability again:
    \[
        P(F \cap E) = P(E|F)P(F)
    \]
    Substituting this back into our expression for $P(F|E)$, we get:
    \[
        P(F|E) = \frac{P(E|F)P(F)}{P(E)}
    \]
\end{proof}

\begin{eg}
    Using the same example of the light bulbs, we want to compute the probability that a randomly chosen bulb was produced by factory A given that it works for over $5000$ hours. We can use Bayes' theorem with $F = B_A$ and $E = \text{works}$:
    \[
        P(B_A|\text{works}) = \frac{P(\text{works}|B_A)P(B_A)}{P(\text{works})}
    \]
    Substituting the values we computed earlier, we get:
    \[
        P(B_A|\text{works}) = \frac{(0.99)(0.60)}{0.974} \approx 0.609
    \]
\end{eg}

\subsection{Random Variables and Probability Distributions}
\begin{definition}[Random Variable]
    A random variable is a function $X: \Omega \to \mathbb{R}$ that assigns a real number to each outcome in the sample space.
\end{definition}

\begin{definition}[Probability Distribution]
    The probability distribution $P_X(x)$ of a random variable $X$ is the probability that $X$ takes the value $x$ i.e. the probability of the event:
    \[
        E = \{\omega \in \Omega : X(\omega) = x\}
    \]
    Hence:
    \[
        P_X(x) = \sum_{\omega \in E} P(\omega)
    \]
\end{definition}
Note that it is often written as:
\[
    P(X = x) = P_X(x)
\]

\begin{eg}
    Let's consider the game of rolling a fair six-sided die and if the outcome is $6$, you receive $10\$$ otherwise you pay $1\$$. We can define a random variable $X$ that represents the amount of money you win (or lose) in this game. The sample space is:
    \[
        \Omega = \{1, 2, 3, 4, 5, 6\}
    \]
    and the random variable $X$ is defined as:
    \[
        X(\omega) = \begin{cases}
            -1 & \text{if } \omega \in \{1, 2, 3, 4, 5\} \\
            10 & \text{if } \omega = 6
        \end{cases}
    \]
    The probability distribution of $X$ is:
    \[
        P_X(-1) = P(X = -1) = \frac{5}{6}, \quad P_X(10) = P(X = 10) = \frac{1}{6}
    \]
\end{eg}
Instead of computing the probability that $X$ takes a specific value, we can compute, for example, the probability that $X \in [a,b]$ for some $a,b \in \mathbb{R}$. This can be done in two ways:
\begin{citemize}
    \item using $P$:
    \[
        \sum_{\omega \in G} P(\omega)
    \]
    where $G = \{\omega \in \Omega : X(\omega) \in [a,b]\}$.
    \item using $P_X$:
    \[
        \sum_{x \in [a,b]} P_X(x)
    \]
\end{citemize}

\begin{eg}
    Lisa rolls two dices and annouces the sum $L$. We want to compute the distribution of $L$. The sample space is:
    \[
        \Omega = \{(i,j) : i,j \in \{1,2,3,4,5,6\}\}
    \]
    and the random variable $L$ is defined as:
    \[
        L(i,j) = i + j
    \]
    The probability distribution of $L$ is:
    \begin{align*}
        P_L(2) &= P(L = 2) = P(\{(1,1)\}) = \frac{1}{36} \\
        P_L(3) &= P(L = 3) = P(\{(1,2), (2,1)\}) = \frac{2}{36} \\
        P_L(4) &= P(L = 4) = P(\{(1,3), (2,2), (3,1)\}) = \frac{3}{36} \\
        &\vdots \\
        P_L(7) &= P(L = 7) = P(\{(1,6), (2,5), (3,4), (4,3), (5,2), (6,1)\}) = \frac{6}{36} \\
        &\vdots \\
        P_L(12) &= P(L = 12) = P(\{(6,6)\}) = \frac{1}{36}
    \end{align*}
\end{eg}

% TODO: add more explanations here when possible (to get slowly twoards two randoms variables and their joint distribution, marginal distribution, etc.)
% TODO: maybe also add examples of why is it useful to have two random variables defined on the same sample space (e.g. the sum of two dice and the first die, etc.)
\begin{definition}[Joint Distribution]
    Let $X$ and $Y$ be two random variables defined on the same sample space $\Omega$. The joint distribution of $X$ and $Y$ is the probability distribution of the pair $(X,Y)$, which assigns a probability to each possible pair of values $(x,y)$ that $X$ and $Y$ can take. Formally, it is defined as:
    \[
        P_{X,Y}(x,y) = P(X = x, Y = y) = P(\{\omega \in \Omega : X(\omega) = x \text{ and } Y(\omega) = y\})
    \]
\end{definition}

\begin{definition}[Marginal Distribution]
    Let $X$ and $Y$ be two random variables defined on the same sample space $\Omega$. The marginal distribution of $X$ is the probability distribution of $X$ obtained by summing over all possible values of $Y$. Formally, it is defined as:
    \[
        P_X(x) = \sum_{y} P_{X,Y}(x,y)
    \]
    where $P_{X,Y}(x,y)$ is the joint distribution of $X$ and $Y$.
\end{definition}

\begin{eg}
    Continuing with the example of Lisa rolling two dice, where $L = L_1 L_2$ is the sum written as two digit number. The alphabet of $L$ is:
    \[
        L = \{02, 03, 04, 05, 06, 07, 08, 09, 10, 11, 12\}
    \]
    Thus, the alphabet of $L_1$ is $\{0,1\}$ and the alphabet of $L_2$ is $\{0,1,2,3,4,5,6,7,8,9\}$. Let's compute the probability $P_{L_1}$, knowing $P_L$. We marginalize:
    \[
        P_{L_1}(1) = \sum_{x} P_{L_1, L_2}(1,x) = P_L(10) + P_L(11) + P_L(12) = \frac{3}{36} + \frac{2}{36} + \frac{1}{36} = \frac{6}{36} = \frac{1}{6}
    \]
\end{eg}

\subsection{Expectation and Independence of Random Variables}
\begin{definition}[Expected Value]
    The expected value (or expectation) of a random variable $X$ is defined as:
    \[
        \mathbb{E}[X] = \sum_{\omega}X(\omega)P(\omega) = \sum_{x} x P_X(x)
    \]
    where the sum is taken either over all outcomes $\omega$ in the sample space or over all values $x$ in the alphabet of $X$.
\end{definition}
Let's show that these two definitions are equivalent. We have:
\begin{align*}
    \sum_{\omega}X(\omega)P(\omega) &= \sum_{x} \sum_{\omega : X(\omega) = x} X(\omega)P(\omega) \\
    &= \sum_{x} x \sum_{\omega : X(\omega) = x} P(\omega) \\
    &= \sum_{x} x P_X(x)
\end{align*}

\begin{eg}
    Let's consider the game of rolling a fair six-sided die and if the outcome is $6$, you receive $10\$$ otherwise you pay $1\$$. We can define a random variable $X$ that represents the amount of money you win (or lose) in this game. The expected value of $X$ is:
    \[
        \mathbb{E}[X] = (-1) \cdot \frac{5}{6} + 10 \cdot \frac{1}{6} = -\frac{5}{6} + \frac{10}{6} = \frac{5}{6} \approx 0.83
    \]
    This means that, on average, you can expect to win about $0.83\$$ per game in the long run.
\end{eg}

\begin{theorem}[Linearity of Expectation]
    Let $X_1, X_2, \dots, X_n$ be $n$ random variables and let $\alpha_1, \alpha_2, \dots, \alpha_n$ be real numbers. Then we have:
    \[
        \mathbb{E}[\alpha_1 X_1 + \alpha_2 X_2 + \dots + \alpha_n X_n] = \alpha_1 \mathbb{E}[X_1] + \alpha_2 \mathbb{E}[X_2] + \dots + \alpha_n \mathbb{E}[X_n]
    \]
\end{theorem}
\begin{proof}
    We can write the expected value of the linear combination of the random variables as:
    \[
        \mathbb{E}[\alpha_1 X_1 + \alpha_2 X_2 + \dots + \alpha_n X_n] = \sum_{\omega} (\alpha_1 X_1(\omega) + \alpha_2 X_2(\omega) + \dots + \alpha_n X_n(\omega)) P(\omega)
    \]
    We can distribute the sum over the linear combination:
    \[
        = \alpha_1 \sum_{\omega} X_1(\omega) P(\omega) + \alpha_2 \sum_{\omega} X_2(\omega) P(\omega) + \dots + \alpha_n \sum_{\omega} X_n(\omega) P(\omega)
    \]
    Recognizing that each sum is the expected value of the corresponding random variable, we get:
    \[
        = \alpha_1 \mathbb{E}[X_1] + \alpha_2 \mathbb{E}[X_2] + \dots + \alpha_n \mathbb{E}[X_n]
    \]
\end{proof}

\begin{definition}[Independence of Random Variables]
    Two random variables $X$ and $Y$ are said to be independent if for all values $x$ and $y$, we have:
    \[
        P_{X,Y}(x,y) = P_X(x)P_Y(y)
    \]
\end{definition}
This definition can be generalized to $n$ random variables $X_1, X_2, \dots, X_n$ by requiring that for all values $x_1, x_2, \dots, x_n$, we have:
\[
    P_{X_1, X_2, \dots, X_n}(x_1, x_2, \dots, x_n) = P_{X_1}(x_1)P_{X_2}(x_2) \cdots P_{X_n}(x_n)
\]

\begin{definition}[Conditional Distribution]
    Let $X$ and $Y$ be two random variables defined on the same sample space $\Omega$. The conditional distribution of $X$ given $Y = y$ is the probability distribution of $X$ restricted to the event $\{Y = y\}$. Formally, it is defined as:
    \[
        P_{X|Y}(x|y) = P(X = x | Y = y) = \frac{P(X = x, Y = y)}{P(Y = y)}
    \]
    provided that $P(Y = y) \neq 0$.
\end{definition}
Note that if $x$ cannot happen (i.e. $P(X = x) = 0$), then the conditional distribution $P_{X|Y}(x|y)$ is not defined, since we cannot condition on an event that cannot happen.

\begin{eg}
    Let's show that the expected value of a product of two independent random variables is NOT always the product of the expected values. Consider $X = Y \in \{-1, 1\}$, uniformly distributed. We have:
    \[
        \mathbb{E}[XY] = 1 \quad \text{and} \quad \mathbb{E}[X]\mathbb{E}[Y] = 0
    \]
    However, if $X$ and $Y$ are independent, then we have:
    \[
        \mathbb{E}[XY] = \mathbb{E}[X]\mathbb{E}[Y]
    \]
\end{eg}

\section{Sources and Entropy}
In the digital world, we communicate by revealing the value of a sequence of variables that we call (information) symbols. For example, the $i$-th symbol might represent:
\begin{citemize}
    \item the intensity of the pixel at position $i$ of a black and white image,
    \item your score in your $i$-th exam,
    \item the $i$-th letter of a text,
\end{citemize}

\subsection{Primitive Definitions of Entropy}
In an article published in 1928 by Hartley (Bell Labs), he wrote that the value of a symbol provides information only if the symbol is a (non-constant) random variable. He the gives a tentative answer to the question, how much information is carried by a symbol such as $S$:
\begin{citemize}
    \item Suppose that $S \in \mathcal{A}$ is a symbol that can take $|\mathcal{A}|$ different values.
    \item The amount of information conveyed by $n$ such symbols should be $n$ times the information conveyed by one symbol.
    \item There are $|\mathcal{A}|^n$ possible values for the $n$ symbols.
    \item This suggest that $\log|\mathcal{A}|^n = n \log |\mathcal{A}|$ is the appropriate measure for information, where we are free to choose the base for the logarithm.
\end{citemize}

\begin{eg}
    In a village that has $8$ telephones, we can assign a different three-digit binary number, such as $001$, to each phone. \\
    Hence it takes $3$ bits of information to identify a phone. Mathematically, the phones are represented by a uniformly distributed random variable $S \in \mathcal{A} = \{1,2, \ldots, 8\}$.
\end{eg}

\begin{eg}
    The world population in 2024 is estimated to be $8.1$ billion. \\
    Hence it takes $\log_2(8.1 \cdot 10^9) = 32.9$ bits of information to identify a person. A person is represented by a uniformly distributed random variable $S \in \mathcal{A} = \{1, 2, \ldots, 8.1 \cdot 10^9\}$.
\end{eg}
Another example shows that something is not right with Hartley's measure of information:

\begin{eg}
    Suppose that $S_n = \{\text{sunny}, \text{rainy}\}$ is the weather prognosis for day $n + 1$, revealed on day $n$. Suppose that $S_n = \text{rainy}$ has a probability of $\frac{5}{365}$. \\
    Hartley's measure assings $\log_2(2) = 1$ bits of information to both, $S_n = \text{sunny}$ and $S_n = \text{rainy}$.
\end{eg}
It seems intuitively obvious that the amount of information provided by $S_n = \text{rainy}$ is much higher than that provided by $S_n = \text{sunny}$ since the likelyhood of rain might be very low. In an article that appeared in 1948, Shannon fixes the problem by defining the notion of uncertainty or entropy $H(S)$ associated to a discrete random variable $S$.

\begin{definition}[Entropy]
    The entropy (or uncertainty) is defined as follows:
    \[
        H_b(S) = - \sum_{s \in supp(P_S)} P_S(s) \log_b(P_S(s))
    \]
    where $supp(P_S) = \{s: P_S(s) > 0\}$.
\end{definition}
The support ($supp$) is to remove the event with a probability of $0$, since $\log$ of $0$ is not defined. Another way to see it more simply is to defined $P_S(s)\log_b(P_S(s)) = 0$ when $P_S(s) = 0$. This convention allows us to simplify the notation to:
\[
    H_b(S) = - \sum_{s \in \mathcal{A}} P_S(s) \log_b(P_S(s))
\]
Note that typically $b = 2$ and in this case the unit is the bit.

\begin{eg}
    When $P_S$ is the uniform distribution over the alphabet $\mathcal{A}$, $P_S(s) = \frac{1}{|\mathcal{A}|}$ and we have:
    \[
        - \log (P_S(s)) = \log |\mathcal{A}|
    \]
    which is constant. In this case, we have:
    \[
        H(s) = \mathbb{E}[\log |\mathcal{A}|] = \log |\mathcal{A}|
    \]
    which is Hartley's information measure.
\end{eg}
Hence Shannon's entropy equals Hartley's measure of information if and only if the random variable has uniform distribution.

\begin{eg}
    A sequence of $4$ decimal digits representing the number to open Anne's lock can be seen as the output of a source $S \in \mathcal{A} = \{1, 2, \ldots, 9998, 9999\}$. If Anne picks the number at random, then we have:
    \[
        P(s) = \frac{1}{10^4} \quad \text{for all $4$-digit numbers}
    \]
    Since the distribution is uniform over the alphabet, we have:
    \[
        H(S) = \log_2 |\mathcal{A}| = \log_2 10^4 \approx 13.3 \text{ bits}
    \]
\end{eg}

\begin{definition}[Binary Entropy Function]
    When $|\mathcal{A}| = 2$, $P_S$ has only two possible values, say $p$ and $(1-p)$. The corresponding entropy is $H(S) = h(p)$ where:
    \[
        h(p) = -p \log_2 p - (1-p)\log_2 (1-p)
    \]
    $h(p)$ is called the binary entropy function.
\end{definition}
Note that:
\begin{citemize}
    \item For $p = 0$ and for $p = 1$, $h(p) = 0$.
    \item For $p = \frac{1}{2}$, $h(p) = 1$.
    \item For $p \in \{0.000001, 0.99999\}$, $H(S) \approx 0.001$.
\end{citemize}
% TODO: maybe add graph (see slide 76)

\begin{eg}
    Let $S_n$ be, like in the previous example, the weather forecast. The probability of $S_n$ are $p = \frac{5}{365}$ and $(1-p) = \frac{360}{365}$. We have:
    \[
        H_2(S_n) = h(\frac{5}{365}) = - \frac{5}{365} \log_2 \frac{5}{365} - \frac{360}{365} \log_2 \frac{360}{365} \approx 0.072 \text{ bits}
    \]
\end{eg}

\begin{eg}
    Taking the previous example, with Anne's lock, the answer to the question "Is $6767$ the lock's number?" can only be $S \in \{\text{YES}, \text{NO}\}$ which is a binary random variable with $P_S(\text{YES}) = \frac{1}{10^4}$. Hence:
    \[
        H(S) = h\left(\frac{1}{10^4}\right) \approx 0.001 \text{ bits}
    \]
\end{eg}
This can be also view as, the answer to the question gives us almost no more information since there is still $9999$ possible combinations.

\begin{definition}[Crazy Distribution]
    The Crazy Distribution is, for a given $a \in \mathcal{A}$:
    \[
        P_S (a) = 1 \quad \text{and} \quad P_S(s) = 0 \text{ for all $s \neq a$}
    \]
\end{definition}
Remark that often, we use the crazy distribution or the uniform distribution as edge cases.

\subsection{Entropy Bounds and Information-Theory Inequality}

\begin{theorem}[IT Inequality]
    For a positive real number $r$, we have:
    \[
        \log_b r \leq (r - 1) \log_b(e)
    \]
    with equality if and only if $r = 1$.
\end{theorem}
\begin{proof}
    Since:
    \[
        \log_b (r) = \ln (r) \log_b(e)
    \]
    we can simply show that $\ln(r) \leq (r - 1)$ with equality if $r = 1$. Graphically this is easy to prove since:
    \begin{citemize}
        \item the function $\ln r$ and $r -1$ coincide at $r = 1$,
        \item the function $r - 1$ has a slope $1$ throughout,
        \item $\frac{d}{dr} \ln(r) = \frac{1}{r} < 1$, for $r > 1$,
        \item $\frac{d}{dr} \ln(r) = \frac{1}{r} > 1$, for $r < 1$.
    \end{citemize}
\end{proof}
% TODO: add graph of ln(r) and (r - 1)

\begin{theorem}[Entropy Bounds]
    The entropy of a discrete random variable $S \in |\mathcal{A}|$ satisfies:
    \[
        0 \leq H_b(S) \leq \log_b |\mathcal{A}|
    \]
    with equality on the left side if and only if $P_S(s) = 1$ for some $s$ (Crazy Distribution), and with equality on the right side if and only if $P_S(s) = \frac{1}{|\mathcal{A}|}$ for all $s$.
\end{theorem}
\begin{proof}
    Let's prove the left side first. Recall that:
    \[
        H(S) = \sum_{s \in |\mathcal{A}|} - P_S(s) \log_b (P_S(s))
    \]
    and we have:
    \[
        -P_S(s) \log_b (P_S(s)) = \begin{cases*}
            0, & if $P_S(s) \in \{0,1\}$ \\
            > 0, & if $0 < P_S(s) < 1$
        \end{cases*}
    \]
    Thus we have that $H(S) \geq 0$, with equality if and only if $P_S(s) \in \{0,1\}$ for all $s \in |\mathcal{A}|$. \\
    To prove the right side, we use a trick that often works in inequalities involving entropies. To prove $A \leq B$ we prove that $A - B \leq 0$ by using the IT Inequality. Thus we have:
    \begin{align*}
        H(S) - \log |\mathcal{A}| &= \mathbb{E}\left[- \log P_S(S)\right] - \log |\mathcal{A}| \\
        &= \mathbb{E}\left[- \log P_S(S) - \log |\mathcal{A}|\right] \\
        &= \mathbb{E}\left[\log \frac{1}{P_S(s) |\mathcal{A}|}\right] \\
        &= \sum_{s \in \mathcal{A}} P_S(s) \left[\log \frac{1}{P_S(s) |\mathcal{A}|}\right] \\
        &\leq \sum_{s \in \mathcal{A}} P_S(s) \left[\frac{1}{P_S(s) |\mathcal{A}|} - 1\right] \log(e) \\
        &= \log(e) \sum_{s \in \mathcal{A}} \left[\frac{1}{|\mathcal{A}|} - P_S(s)\right] \\
        &= \log(e)(1-1) = 0
    \end{align*}
    with equality if and only if $P_S(s) |\mathcal{A}| = 1$ for all $s$.
\end{proof}

\begin{eg}
    Let $S$ be Anne's lock number. Its entropy is maximized if Anne chooses at random over all $10^4$ possibilities. In this case, and only in this case:
    \[
        H(S) = \log |\mathcal{A}| = \log 10^4 \approx 13.3 \text{ bits}
    \]
\end{eg}

\subsection{Entropy of Multiple Random Variables}
The formula for the entropy of a random variable can be extended to any number of random variables. If $X$ and $Y$ are two discrete random variable, with (joint) probability distribution $P_{X,Y}$ then:
\[
    H(X, Y) = \mathbb{E}\left[- \log P_{X, Y}(X,Y)\right]
\]
which means:
\[
    H(X,Y) = - \sum_{(x, y) \in \mathcal{X} \times \mathcal{Y}} P_{X, Y}(x, y) \log P_{X, Y}(x,y)
\]

\begin{eg}
    Let $P_{X,Y}$ be given by:
    \[
        P_{X,Y}(0,0) = \frac{1}{8}, \quad P_{X,Y}(0,1) = \frac{3}{8}, \quad P_{X,Y}(1,0) = \frac{1}{4}, \quad P_{X,Y}(1,1) = \frac{1}{4}
    \]
    Thus we can compute the entropy:
    \[
        H(X,Y) = -\frac{1}{8} \log_2\frac{1}{8} - \frac{3}{8} \log_2 \frac{3}{8} - \frac{1}{4} \log_2 \frac{1}{4} - \frac{1}{4} - \log_2 \frac{1}{4}
    \]
\end{eg}
We are mainly intrested in sources that emit a large number of random variables. The sequence of random variables can be:
\begin{citemize}
    \item Finite, like in $(S_1, \ldots, S_n)$ (random vector).
    \item Infinite, like in $S_1, S_2, \ldots$ (random sequence, random process), also denoted by $\{S_i\}_{i = 1}^\infty$.
    \item Sometimes it is convenient to consider $\ldots, S_{-1}, S_0, S_1, \ldots$, also denoted by $\{S_i\}$.
\end{citemize}

\begin{eg}
    A sequence of coin flips can be seen as the output of a source $S_1, S_2, \ldots, S_n$ with $S_i \in \mathcal{A} = \{\text{H}, \text{T}\}$, where H stands for head, and T for tail, for $i = 1,\ldots, n$. If the coin is fair, all sequence are equally likely:
    \[
        P(s_1, s_2, \ldots, s_n) = \prod_{i}P(s_i) = \frac{1}{2^n} \text{ for all } (s_1, s_2, \ldots, s_n) \in \mathcal{A}^n
    \]
\end{eg}

\section{The Fundamental Compression Theorem}
Our first concrete application of the notion of entropy is in source coding (i.e. data compression). Source coding is the art of representing the output of a source (i.e. a sequence of random variables) using as few bits as possible. 

\begin{definition}[Source]
    The source is specified by the source alphabet $\mathcal{A}$ and by the distribution statistics.
\end{definition}

\begin{eg}
    The source alphabet is $\mathcal{A} = \{a, \ldots, z, 0, \ldots, 9\}$ and the source symbol are independent and identically distributed (i.i.d.) over $\mathcal{A}$.
\end{eg}

\begin{definition}[Encoder]
    An encoder is specified by:
    \begin{citemize}
        \item the input alphabet $\mathcal{A}$ (the same as the source alphabet),
        \item the output alphabet $\mathcal{D}$ (typically $\mathcal{D} = \{0, 1\}$),
        \item the codebook $\mathcal{C}$ which consists of a finite sequence over $\mathcal{C}$,
        \item by the one-to-one encoding map (encoding function) $\Gamma: \mathcal{A}^k \to \mathcal{C}$, where $k$ is a positive integer. 
    \end{citemize}
\end{definition}
For now, $k = 1$.

\begin{eg}
    The source alphabet is $\mathcal{A} = \{a,b,c,d\}$, with the following encoding maps:
    \[
        \Gamma_O(a) = 00, \quad \Gamma_O(b) = 01, \quad \Gamma_O(c) = 10, \quad \Gamma_O(d) = 11
    \]
    \[
        \Gamma_A(a) = 0, \quad \Gamma_A(b) = 01, \quad \Gamma_A(c) = 10, \quad \Gamma_A(d) = 11
    \]
    \[
        \Gamma_B(a) = 0, \quad \Gamma_B(b) = 10, \quad \Gamma_B(c) = 110, \quad \Gamma_B(d) = 1110
    \]
    \[
        \Gamma_C(a) = 0, \quad \Gamma_C(b) = 01, \quad \Gamma_C(c) = 011, \quad \Gamma_C(d) = 0111
    \]
    Thus for example, if we want to encode the word "dac", we have:
    \[
        \Gamma_O(dac) = 111010, \quad \Gamma_A(dac) = 11010, \quad \Gamma_B(dac) = 11100110, \quad \Gamma_C(dac) = 01110011
    \]
    The source alphabet $\mathcal{A}$, the $k$, the code alphabet $\mathcal{D}$ and the codebook $\mathcal{C}$ are implicitly defined by the encoding maps $\Gamma_O, \Gamma_A, \Gamma_B, \Gamma_C$.
\end{eg}

\subsection{Decodability}
We want to avoid the following situation.
\begin{eg}
    If we use the encoding map $\Gamma_A$ from the previous example, the word "cbaad" is encoded as $10010011$ which can be decoded as "cbaad" but also as "cacad". 
\end{eg}

\begin{definition}[Uniquely Decodable Code]
    The code is uniquely decodable if every concatenation of codewords has a unique parsing into a sequence of codewords.
\end{definition}
Uniquely decodable codes allow us to store (or transmit) a sequence of symbols without ambiguity and without separators between codewords.

\begin{eg}
    The code $\Gamma_A$ is not uniquely decodable since:
    \[
        \Gamma_A(bc) = 0110 = \Gamma_A(ada)
    \]
    the same happens for $0100$.
\end{eg}

\begin{eg}
    The code $\Gamma_B$ is uniquely decodable since every $0$ marks the end of a codeword (i.e. the $0$ plays the role of a separator). The same can be said for $\Gamma_C$ since every codeword starts with $0$.
\end{eg}

\begin{eg}
    The code $\Gamma_O$ is uniquely decodable. A fixed-length code is always uniquely decodable.
\end{eg}

\begin{definition}[Prefix-Free Code]
    If no codeword is a prefix of another codeword, then the code is called prefix-free.
\end{definition}

\begin{eg}
    The codeword $01$ is a prefix of the codeword $011$ hence $\Gamma_C$ is not prefix-free. On the other hand, $10$ is not a prefix of $110$.
\end{eg}
Note that:
\begin{citemize}
    \item A prefix-free code is always uniquely decodable,
    \item A uniquely decodable code is not necessarily prefix-free.
\end{citemize}

\begin{eg}
    The code $\Gamma_C$ is not prefix-free (since $01$ is a prefix of $011$) but it is uniquely decodable since every codeword starts with $0$. It's reverse (i.e. $\Gamma_B$) is prefix-free.
\end{eg}
Remark that a code is uniquely decodable if and only if its reverse is uniquely decodable. A prefix-free code is also called instantaneous code since it can be decoded as soon as the last bit of a codeword is received. 

\subsection{Codes for One Random Variable}
We start by considering codes that encode one single random variable $S \in \mathcal{A}$. To encode a sequence $S_1, S_2, \ldots$ of random variables, we encode one random variable at a time.

\subsection{Tree of a Code}

\begin{eg}
    The codes $\Gamma_O$ and $\Gamma_A$, on the alphabet $\mathcal{A} = \{a,b,c,d\}$, where:
    \[
        \Gamma_O(a) = 00, \quad \Gamma_O(b) = 01, \quad \Gamma_O(c) = 10, \quad \Gamma_O(d) = 11
    \]
    \[
        \Gamma_A(a) = 0, \quad \Gamma_A(b) = 01, \quad \Gamma_A(c) = 10, \quad \Gamma_A(d) = 11
    \]
    They can be represented by the following trees:
    % TODO: make the trees
\end{eg}

\begin{definition}[Decoding Tree]
    A decoding tree can be constructed from a complete tree by keeping only the branches that form a codeword.
\end{definition}

\begin{definition}[Codeword Length]
    The length of a codeword is the number of branches that we need to follow in the decoding tree to reach the codeword.
\end{definition}
% TODO: add example of codeword length for the previous example

\subsection{Kraft-McMillan}

\begin{theorem}[Kraft Inequality]
    If a $D$-ary code is uniquely decodable then its codeword lengths $l_1, \ldots, l_M$ satisfy:
    \[
        D^{-l_1} + \ldots + D^{-l_M} \leq 1
    \]
\end{theorem}
% \begin{proof}
%     We will prove a slightly weaker result, namely that the codeword lengths of prefix-free codes satisfy the Kraft inequality.
% \end{proof}
Note that the Kraft inequality is a necessary condition for unique decodability but not a sufficient condition.

\begin{eg}
    For the code $\Gamma_O$, we have:
    \[
        2^{-2} + 2^{-2} + 2^{-2} + 2^{-2} = 1
    \]
    Hence Kraft's inequality is fulfilled with equality.
\end{eg}

\begin{eg}
    For the code $\Gamma_C$, which is given by:
    \[
        \Gamma_C(a) = 0, \quad \Gamma_C(b) = 01, \quad \Gamma_C(c) = 011, \quad \Gamma_C(d) = 0111
    \]
    the codeword lengths are $l_a = 1$, $l_b = 2$, $l_c = 3$ and $l_d = 4$. Hence we have:
    \[
        2^{-1} + 2^{-2} + 2^{-3} + 2^{-4} = \frac{15}{16} = 0.9375 < 1
    \]
    Thus Kraft's inequality is fulfilled but $\Gamma_C$ is not uniquely decodable.
\end{eg}
Note that the contrapositive of Kraft's inequality is a sufficient condition to prove that a code is not uniquely decodable. 

\begin{eg}
    For the code $\Gamma_A$, which is given by:
    \[
        \Gamma_A(a) = 0, \quad \Gamma_A(b) = 01, \quad \Gamma_A(c) = 10, \quad \Gamma_A(d) = 11
    \]
    the codeword lengths are $l_a = 1$, $l_b = 2$, $l_c = 2$ and $l_d = 2$. Hence we have:
    \[
        2^{-1} + 2^{-2} + 2^{-2} + 2^{-2} = \frac{5}{4} = 1.25 > 1
    \]
    Thus Kraft's inequality is not fulfilled and there exists no uniquely decodable code with the same codeword lengths as $\Gamma_A$.
\end{eg}

\begin{theorem}[Kraft-McMillan Theorem]
    For a given set of codeword lengths $l_1, \ldots, l_M$, there exists a uniquely decodable code with these codeword lengths if and only if the Kraft inequality is satisfied i.e. if and only if:
    \[
        D^{-l_1} + \ldots + D^{-l_M} \leq 1
    \]
\end{theorem}

\begin{eg}
    Let $l_1 = 1$, $l_2 = 2$, $l_3 = 3$ and $l_4 = 4$. We have:
    \[
        2^{-1} + 2^{-2} + 2^{-3} + 2^{-4} = \frac{15}{16} < 1
    \]
    Hence there exists a uniquely decodable code with these codeword lengths. For example, we can use the code $\Gamma_C$ from the previous example, which is given by:
    \[
        \Gamma_C(a) = 0, \quad \Gamma_C(b) = 01, \quad \Gamma_C(c) = 011, \quad \Gamma_C(d) = 0111
    \]
\end{eg}

\subsection{Average Codeword Length}
The typical use of a code is the encode a sequence of random variables into the corresponding codeword sequence. We are interested in minimizing the average number of bits per symbol in the codeword sequence.

\begin{definition}[Average Codeword Length]
    Let $l(\Gamma(s))$ be the length of the codeword $\Gamma(s)$ corresponding to the symbol $s$. The average codeword length is defined as:
    \[
        L(S, \Gamma) = \sum_{s \in \mathcal{A}} P_S(s) l(\Gamma(s))
    \]
\end{definition}
The units of $L(S, \Gamma)$ are code symbols, when $D = 2$ the units are bits.

\begin{eg}
    For the following codes:
    \[
        \Gamma_O(a) = 00, \quad \Gamma_O(b) = 01, \quad \Gamma_O(c) = 10, \quad \Gamma_O(d) = 11
    \]
    \[
        \Gamma_B(a) = 0, \quad \Gamma_B(b) = 10, \quad \Gamma_B(c) = 110, \quad \Gamma_B(d) = 1110
    \]
    \[
        \Gamma_B(a) = 1110, \quad \Gamma_B(b) = 110, \quad \Gamma_B(c) = 10, \quad \Gamma_B(d) = 0
    \]
    \[
        \Gamma_{B'}(a) = 0, \quad \Gamma_{B'}(b) = 10, \quad \Gamma_{B'}(c) = 110, \quad \Gamma_{B'}(d) = 1110
    \]
    associated to the alphabet $\mathcal{A} = \{a,b,c,d\}$, with the following probabilities:
    \[
        P_S(a) = 0.05, \quad P_S(b) = 0.05, \quad P_S(c) = 0.1, \quad P_S(d) = 0.8
    \]
    we have:
    \[
        L(S, \Gamma_B) = 0.05 \cdot 1 + 0.05 \cdot 2 + 0.1 \cdot 3 + 0.8 \cdot 4 = 3.65
    \]
    \[
        L(S, \Gamma_{B'}) = 0.05 \cdot 4 + 0.05 \cdot 3 + 0.1 \cdot 2 + 0.8 \cdot 1 = 1.35
    \]
    Thus we have:
    \[
        L(\Gamma_{B'}) < L(\Gamma_O) < L(\Gamma_B)
    \]
\end{eg}
We now want to find a lower bound on the average codeword length. 

\begin{theorem}
    Let $\Gamma: \mathcal{A} \to \mathcal{C}$ be the encoding map of a D-ary code for the random variable $S \in \mathcal{A}$. If the code is uniquely decodable, then the average codeword length is lower bounded by the entropy of $S$, namely:
    \[
        H_D(S) \leq L(S, \Gamma)
    \]
    with equality if and only if $P_S(s) = D^{-l(\Gamma(s))}$ for all $s \in \mathcal{A}$. An equivalent condition is $l(\Gamma(s)) = \log_D \frac{1}{P_S(s)}$.
\end{theorem}
\begin{proof}
    We have:
    \begin{align*}
        H_D(S) - L(S, \Gamma) &= \sum_{i} P_i \log_D \frac{1}{P_i} - \sum_{i} P_i \log_D D^{l_i} \\
        &= \sum_{i} P_i \log_D \frac{1}{P_i D^{l_i}} \\
        &\leq \log_D (e) \sum_{i} P_i \left(\frac{1}{P_i D^{l_i}} - 1\right) & \text{(IT inequality)} \\
        &= \log_D (e) \sum_{i} \left(\frac{1}{D^{l_i}} - P_i\right) \\
        &= \log_D (e) \left(\sum_{i} \frac{1}{D^{l_i}} - 1\right) \\
        &\leq 0 & \text{(Kraft's inequality)}
    \end{align*}
    The first inequalities hold with equality if and only if, for all $i$, we have:
    \[
        P_i = \frac{1}{D^{l_i}}
    \]
    When this is the case, the second inequality holds with equality.
\end{proof}

\begin{eg}
    Using the same codes as in the previous example, we have:
    \[
        H(S) = -2 \cdot 0.05 \log_2 0.05 - 0.1 \log_2 0.1 - 0.8 \log_2 0.8 \approx 1.022 \text{ bits}
    \]
    Hence we have:
    \[
        H(S) < L(S, \Gamma_{B'}) < L(S, \Gamma_O) < L(S, \Gamma_B)
    \]
\end{eg}

\subsection{Shannon-Fanno Codes}
A good observation from the previous theorem is that the right side of both equations are similar:
\[
    L(S, \Gamma) = \sum_{s \in \mathcal{A}} P_S(s) l(\Gamma(s)) \quad \text{and} \quad H_D(S) = \sum_{s \in \mathcal{A}} P_S(s) \log_D \frac{1}{P_S(s)}
\]
and thus are equal if and only if $l(\Gamma(s)) = \log_D \frac{1}{P_S(s)}$ for all $s \in \mathcal{A}$.

\begin{definition}[Shannon-Fano Code]
    A Shannon-Fano code is a code such that for all $s \in \mathcal{A}$, we have:
    \[
        l(\Gamma(s)) = \lceil \log_D \frac{1}{P_S(s)} \rceil
    \]
\end{definition}

\begin{eg}
    Let's build a Shannon-Fano code for the random variable $S \in \mathcal{A} = \{a,b,c,d\}$ with the following probabilities:
    \[
        P_S(a) = 0.05, \quad P_S(b) = 0.05, \quad P_S(c) = 0.1, \quad P_S(d) = 0.8
    \]
    We have:
    \[
        l(\Gamma(a)) = \lceil \log_2 \frac{1}{0.05} \rceil = 5, \quad l(\Gamma(b)) = \lceil \log_2 \frac{1}{0.05} \rceil = 5
    \]
    \[
        l(\Gamma(c)) = \lceil \log_2 \frac{1}{0.1} \rceil = 4, \quad l(\Gamma(d)) = \lceil \log_2 \frac{1}{0.8} \rceil = 1
    \]
    Hence we can use the following code:
    \[
        \Gamma(a) = 00000, \quad \Gamma(b) = 00001, \quad \Gamma(c) = 0001, \quad \Gamma(d) = 1
    \]
\end{eg}
Note that the Shannon-Fano code is not necessarily optimal, since it does not necessarily minimize the average codeword length.

\begin{theorem}
    The average codeword length of a Shannon-Fano code for the random variable $S$ furfills:
    \[
        H_D(S) \leq L(S, \Gamma_{SF}) < H_D(S) + 1
    \]
\end{theorem}
\begin{proof}
    We will only prove the upper bound, since the lower bound is trivial. We have that:
    \[
        l_i = \lceil \log_D \frac{1}{P_i} \rceil < \log_D \frac{1}{P_i} + 1
    \]
    Hence we have:
    \begin{align*}
        L(S, \Gamma_{SF}) &= \sum_{i} P_i l_i \\
        &< \sum_{i} P_i \left(\log_D \frac{1}{P_i} + 1\right) \\
        &= H_D(S) + 1
    \end{align*}
\end{proof}

\subsection{Huffman Codes}
Unlike Shannon-Fano codes, Huffman codes are optimal, meaning that they minimize the average codeword length. To build a Huffman code, we start